\documentclass{article}
\usepackage[UTF8]{ctex}
\usepackage{amsmath,amssymb,amsfonts}
\usepackage{enumitem}
\usepackage{geometry}
\geometry{a4paper, margin=1in}
% === 错题本排版优化（不参与题目标识） ===
\setlength{\parindent}{0pt}
\setlength{\parskip}{0.35em plus 0.10em minus 0.08em}
\setlength{\emergencystretch}{3em}
\setlist{itemsep=0.25em, topsep=0.40em, parsep=0pt}
\allowdisplaybreaks[2]
% === 错题本排版优化结束 ===
\begin{document}
\section*{南京航空航天大学2024年高等代数真题}
\section*{92．南京航空航天大学2024年高等代数试卷}

\subsection{题目1}
1．给定矩阵

\[
A=\left(\begin{array}{lll}
1 & -3 & 4 \\
4 & -7 & 8 \\
6 & -7 & 7
\end{array}\right)
\]

（1）求 \(A\) 的特征值和最小多项式。\\
（2）求 \(A\) 的初等因子和 Jordan 标准形．

\subsection{题目2}
2．设 \(V_{1}\) 是由

\[
\alpha_{1}=(1,1, a)^{\mathrm{T}}, \alpha_{2}=(-2, a, 4)^{\mathrm{T}}, \alpha_{3}=(-2, a,-2)^{\mathrm{T}}
\]

生成的 \(\mathbb{R}^{3}\) 的子空间，\(V_{2}\) 是由

\[
\beta_{1}=(1,1, a)^{\mathrm{T}}, \beta_{2}=(1, a, 1)^{\mathrm{T}}, \beta_{3}=(a, 1,1)^{\mathrm{T}}
\]

生成的 \(\mathbb{R}^{3}\) 的子空间。

\begin{itemize}
  \item[（1）] 若 \(V_{1} \neq V_{2}\) ，求 \(a\) 的范围．
  \item[（2）] 当 \(a=2\) 时，求 \(\operatorname{dim}\left(V_{1} \cap V_{2}\right)\) ．
\end{itemize}

\subsection{题目3}
3．设 \(\sigma\) 为 \(\mathbb{R}^{3}\) 上的线性变换，\(\varepsilon_{1}=(1,1,0)^{\mathrm{T}}, \varepsilon_{2}=(0,1,1)^{\mathrm{T}}, \varepsilon_{3}=(1,1,1)^{\mathrm{T}}\) 为 \(\mathbb{R}^{3}\) 的一组基，\(\sigma\) 在基\(\varepsilon_{1}, \varepsilon_{2}, \varepsilon_{3}\) 下的矩阵为 \(A\) ，且 \(\sigma\left(\varepsilon_{1}\right)=(1,0,0)^{\mathrm{T}}, \sigma\left(\varepsilon_{2}\right)=(0,1,1)^{\mathrm{T}}, \sigma\left(\varepsilon_{3}\right)=(1,1+a, a)^{\mathrm{T}}\) 。

\begin{itemize}
  \item[（1）] 若 \(\sigma\) 可对角化，求 \(a\) 的值．
  \item[（2）] 当 \(a=2\) 时，求 \(A^{2023}\) 。
\end{itemize}

\subsection{题目4}
4．设 \(\alpha, \beta\) 为 \(\mathbb{R}^{3}\) 中的单位列向量且相互正交，实二次型 \(f(X)=2\left(\alpha^{\mathrm{T}} X\right)^{2}+\left(\beta^{\mathrm{T}} X\right)^{2}\) 的矩阵为 \(A\) ，其中 \(X=\left(x_{1}, x_{2}, x_{3}\right)^{\mathrm{T}}\) 。证明：

\begin{itemize}
  \item[（1）] 存在正交矩阵 \(U\) ，使得 \(U^{\mathrm{T}} A U\) 为对角形 \(\operatorname{diag}\{2,1,0\}\) 。
  \item[（2）] 是否存在唯一的半正定矩阵 \(S\) 使得 \(A=S^{2}\) ？请说明理由。
\end{itemize}

\subsection{题目5}
5．设 \(A, B\) 是两个 \(n\) 阶矩阵，且 \(A B=4 B+3 A-10 E\) 。证明：

\begin{itemize}
  \item[（1）] \(\lambda=4\) 不是 \(A\) 的特征值，且 \(A B=B A\) 。
  \item[（2）] 若 \(A\) 相似于对角矩阵，则存在可逆矩阵 \(P\) ，使得 \(P^{-1} A P, P^{-1} B P\) 均为对角矩阵。
\end{itemize}

\subsection{题目6}
6．设 \(A, B, C, D\) 都是 \(n\) 阶矩阵且 \(A B=B A\) 。证明：

\begin{itemize}
  \item[（1）] 若 \(r(A)=n\) ，则 \(r\left(\begin{array}{cc}A & B \\ C & D\end{array}\right)=n+r(D A-C B)\) 。
  \item[（2）] \(\left|\begin{array}{ll}A & B \\ C & D\end{array}\right|=|D A-C B|\) 。
\end{itemize}

\subsection{题目7}
7．设复方阵 \(A=\left(a_{i j}\right)_{n \times n}\) ，记 \(\bar{A}=\left(\bar{a}_{i j}\right)_{n \times n}\)（ \(\bar{a}_{i j}\) 为 \(a_{i j}\) 的共轭复数）。如果复方阵 \(A\) 满足 \((\bar{A})^{\mathrm{T}}=A\) ，称 \(A\) 为 Hermite 矩阵。证明：

\begin{itemize}
  \item[（1）] 若复方阵 \(A\) 为 Hermite 矩阵，则 \(A\) 的特征值均为实数。
  \item[（2）] 若复方阵 \(A\) 满足 \(A(\bar{A})^{\mathrm{T}}=A^{2}\) ，则 \(A\) 为 Hermite 矩阵。
\end{itemize}

\subsection{题目8}
8．设 \(V\) 是数域 \(P\) 上的 \(n\) 维线性空间，\(\sigma\) 是 \(V\) 上的线性变换，多项式 \(f(x), g(x) \in P[x]\) ，且 \(h(x)=f(x) g(x),(f(x), g(x))=1\) ，记 \(\operatorname{Ker} \sigma=\{\alpha \mid \alpha \in V, \sigma(\alpha)=0\}\) ．证明：

\begin{itemize}
  \item[（1）] \(\operatorname{Ker} h(\sigma)=\operatorname{Ker} f(\sigma) \oplus \operatorname{Ker} g(\sigma)\) ．
  \item[（2）] \(\operatorname{dim} \operatorname{Ker} \sigma+\operatorname{dim} \operatorname{Ker} \sigma^{3} \leq 2 \operatorname{dim} \operatorname{Ker} \sigma^{2}\) ．
\end{itemize}
\end{document}
